
\subsection*{Biodiversity metrics}

We modeled two biodiversity metrics, the geometric mean abundance (GMA) \cite{santini2017}, an alpha diversity metric, and the Bray-Curtis (BC) similarity \cite{bray1957}, a beta diversity metrics. Combined, they provide insights into the relative level of community richness, abundance, evenness, and compositional similarity between areas, and are suitable for detecting changes in biodiversity \cite{santini2017}. If we let $a_{si}$ represent the population abundance of species $s$ at sampling site $i$, and let $S_i$ represent the total number of species at the site, the GMA of that site can be calculated as
\vspace{0.25\baselineskip}
\begin{equation*}
    y_i = \exp\left(\frac{\sum_{s=1}^{S_i} \ln a_{si}}{S_i}\right).
\end{equation*}

The ln transformation dampens the contribution of highly abundant (perhaps generalist or opportunistic) species. For a given total site abundance, higher species richness and evenness results in a greater GMA. For beta diversity, the BC similarity between two sampling sites $i$ and $j$ is given by
\vspace{0.25\baselineskip}
\begin{equation*}
    y_{ij} = \frac{2\sum_{s=1}^S \min(a_{si},a_{sj})}{\sum_{s=1}^S \big(a_{si}+a_{sj}\big)},
\end{equation*}

where $S$ represents the total number of species found at any of the sites. This is equivalent to an abundance extension of the Sørensen index. The index takes a value of 0 if there are no shared species between sites $i$ and $j$, and a value of 1 if the species and their abundances are identical between the two sites (which is highly unlikely). Since the focus of this study is on models used to estimate biodiversity intactness, the BC similarity is calculated between reference sites, consisting of minimally used primary vegetation sites, and all other sites within each study. This follows the approach used in the BII\cite{depalma2021} and is similar to the implementation of MSA in GLOBIO \cite{schipper2020}.

\vspace{0.25\baselineskip}
One challenge with biodiversity meta-databases is that the abundance distributions from different source studies will be greatly influence by taxonomic scope, biogeographic area and sampling method. Furthermore, while abundance is often expressed as a count of individuals, it can also be a proportion or density. The most viable approach we found, following the BII \cite{depalma2021}, was to normalize the GMA values within each study to a common 0--1 scale, by dividing them by the respective within-study maxima:
\begin{equation*}
    y_i^{\text{norm}} = \frac{y_i}{\max(y_1,\ldots,y_I)}.
\end{equation*}

The BC index is already expressed on a 0--1 scale, so no further processing was needed. For simplicity, we let $y_i$ denote the normalized values from hereon. Although the sampling method is consistent between sampling sites within a given source study, sampling effort between sites can vary, so effort-adjusted abundance numbers (already provided in the PREDICTS data) were used in all analyses. Since both diversity metrics require abundance data, we filtered out studies that only recorded presences and absences. To mitigate the impact of extreme abundance values, we identified outlier locations using the interquartile range (IQR) method and removed site-level observations where $y_i > 1.5 \,\text{IQR}$, within each study.

The distributions of the alpha diversity (GMA) and beta diversity (BC similarity) metrics are shown in \autoref{fig:response_data_distr}. They are non-symmetric and quite heavily skewed to the right. The shapes are similar, but the GMA distributions have greater means than the BC indices, since a compositional similarity metric to a greater extent reflects natural species turnover across the landscape. There is an inflation of zeros in both distributions, and an inflation of ones in the GMA data. A site-level abundance of zero can either be the result of true absence of the surveyed species, or the result of imperfect detection, a well-known issue in biodiversity data. In the best case, noise that arises from imperfect detection is randomly distributed across studies and sites, but we acknowledge that there could be more systematic detection biases in the data. In addition, studies that only surveyed a small number of species are more likely to have zeros at the site level. However, there is no feasible way to construct a separate detection probability model \cite{dorazio2016,wu2015} for such a heterogeneous dataset using the data that we have available. The inflation of zeros in the BC distribution is the result of a mix of observed zero abundances and complete dissimilarity in species composition between sites. The inflation of ones in the GMA data is an artifact of the normalization procedure described above, since every study will contribute a one (its maximum abundance site) to the overall data pool.

In the main results, we used data from all studies except for what was subject to the filters described above (such as requiring abundance data). It should be noted, however, that one issue is that many source studies contain very few sites: 18.8\% contain 5 sites or less, and 60.9\% contain 25 or fewer (see \autoref{fig:sites_per_study}). There is a risk that such small studies contribute more noise than signal to the overall pool of data used to train the models. Informally, this is because there are not enough observations from the context of a given study to reliably relate its species observations to different human pressures and environmental conditions. Through inspection of study-level histograms, we found that ideally 25--50 sites per study would be required to produce somewhat continuous distributions of data at the study-level. Some of this is alleviated by pooling data across studies, but some biases can certainly persist.

%This is supported by the fact that a threshold of 25 sites when training and evaluating the BHM resulted in better model performance than no threshold or a 10-site threshold (\hl{Extended Data Fig. X}). We still chose to include all data as the aim of the study was to illustrate the current state of model performance using available data, not optimizing for accuracy.

The biogeographic-taxonomic models (BTMs) used taxonomic groups as part of the hierarchical model structure (see \autoref{fig:hierarchical_groups}). Since some larger studies had sampled species across several such taxonomic groups, sampling sites were consequently split into multiple observations (compare \extfigpanel{fig:response_data_distr}{c,d} to \extfigpanel{fig:response_data_distr}{a,b}). The BTM alpha model had 31,242 observations in total, in comparison to the 25,987 sampling sites in scope, which also equaled the number of observations in the study-block model (SBM). Since the BC similarity dataset contained more than 400,000 unique site pairs after standard filtering, we used sub-sampling to obtain a tractable but representative dataset for model training and evaluation. The sampling was designed to retain all studies with at least one reference site, while preventing large studies, with many sites and reference sites, from becoming too dominant. For each study $s$, we calculated the number of potential pairs $n_{sites} \times n_{ref}$ and a sublinear weight $w_s = \sqrt{n_{sites} \times n_{ref}}$, with $w_{tot}$ denoting the sum of all study weights. An initial number of pairs was allocated to each study based on its actual number of pairs $n_{pairs}$ and an overall target fraction $f  = 0.10$ of the whole dataset. This was weighted by $w_s$ and subject to lower and upper bounds $k_{min}, k_{max}$, such that $k_s = \text{max}(k_{min}, \text{min}(n_{pairs} \times f \times w_s / w_{tot}), k_{max})$. The threshold values were set to $k_{min} = 300$ and $k_{min} = 3,000$, respectively, to increase the contribution of smaller studies and limit the dominance of the largest ones. To balance reference sites and other sites within each study sample, we finally chose $\sqrt{k_{study}}$ non-reference sites and $k_{study} / \sqrt{k_{study}}$ reference site, to get approximately $k_{study}$ sites in total. This resulted in a total of 29,088 site pairs that were consistent across SBM and BTM models.
\vspace{\baselineskip}
